% !TeX spellcheck = en_US
\documentclass[french]{yLectureNote}

\title{Atomistique 2 - PS}
\subtitle{Chimie}
\author{Paulhenry Saux}
\date{\today}
\yLanguage{Français}

\professor{J.Cunny}
\usepackage{graphicx}%----pour mettre des images
\usepackage[utf8]{inputenc}%---encodage
\usepackage{geometry}%---pour modifier les tailles et mettre a4paper
%\usepackage{awesomebox}%---pour les boites d'exercices, de pbq et de croquis ---d\'esactiv\'e pour les TP de PC
\usepackage{tikz}%---pour deiffner + d\'ependance de chemfig
% \usepackage{tabularx}%---pour dimensionner automatiquement les tableaux avec variable X
\usepackage{awesomebox}%---Pour les boites info, danger et autres
\usepackage{menukeys}%---Pour deiffner les touches de Calculatrice
\usepackage{fancyhdr}%---pour les en-t\^ete personnalis\'ees
\usepackage{blindtext}%---pour les liens
\usepackage{hyperref}%---pour les liens (\`a mettre en dernier)
\usepackage{caption}%---pour la francisation de la l\'egende table vers Tableau
\usepackage{pifont}
\usepackage{array}%---pour les tableaux
\usepackage{lipsum}
\usepackage{yFlatTable}
\usepackage{multicol}
\newcommand{\Lim}[1]{\lim\limits_{\substack{#1}}\:}
\renewcommand{\vec}{\overrightarrow}
\newcommand{\N}[0]{\mathbb{N}}
\newcommand{\dd}{\mathrm{d}}
\newcommand{\norm}[1]{||\vec{#1}||}
\newcommand{\fo}{\psi(\vec{r},t)}
\newcommand{\foe}{\psi(\vec{r},t)\*}
\newcommand{\HH}{\hat{H}}
\newcommand{\hb}{\hbar}
\newcommand{\lap}{\nabla^2}
\newcommand{\lapcc}{\frac{\partial^2 }{\partial x^2}+\frac{\partial^2 }{\partial y^2}+\frac{\partial^2 }{\partial z^2}}
\begin{document}
\setcounter{chapter}{3}
\chapter{Symétrie moléculaire}
\section{Opérations}
\criticalInfo{Élément de symétrie et opération}{L'élément de symétrie est l'élément virtuel avec lequel on applique l'opération.}
\subsection{Rotation}
Élément de symétrie : Axe noté \(C_n\) avec \(n\) décrivant l'angle de rotation tel que l'angle vaut \(2\pi/n\).

\begin{tabular}{ll}
C2 & 180°\\
C3 & 120°\\
C4 & 90°\\
C6 & 60°
                     \end{tabular}

Cet axe engendre \(n-1\) opérations de symétries notée \(C_n^m\) avec \(m\) le nombre de fois que l'on applique la rotation\marginTips{On peut diviser m et n par le m\^eme entier. Si on tombe dans les 2 cas sur 2 autres entiers, les 2 opérations sont identiques. Par exemple \(C^2_4 = C^1_2\)}.
\criticalInfo{C22}{\(C_2^2\) n'existe pas car c'est déjà l'opération Identité}
\warningInfo{Axe principal}{Lorsque la molécule possède plusieurs axes de symétrie, l'axe principal est celui avec le \(n\) le plus élevé. Par convention, il est orienté selon \(z\)}
\begin{definition}[Molécule très symétrique]
On dit d'une molécule qu'elle est très symétrique lorsqu'elle possède au moins 2 axes principaux (octaèdre, isocaèdre, icosaèdre)
\end{definition}
On peut définir des axes colinéaires à des axes de grand \(n\) en trouvant leur diviseur. Par exemple \(12 = 2\times 2 \times 3\). Ainsi \(C_{12}\) génère des axes \(C_2,C_4,C_3, C_6\).
\criticalInfo{Axes \(C_2\) multiples}{Dans le cas d'axes \(C_2\) multiples, on note \(C_2\) l'axe engendré par l'axe de plus grand \(n\) si applicable, \(C_2'\) les axes passant par les atomes et \(C_2''\) les axes passant par les milieux des arêtes.}
\criticalInfo{Benzène}{Le Benzène a un axe \(C_6\)}
\subsection{Réflexion}
Élément de symétrie : un plan noté \(\sigma\) qui engendre une seule symétrie. Il n'y a donc pas d'indices. En revanche, on a \(\sigma^m = \sigma\) si \(m\) est impair et \(\sigma^m =E\) si il pair.

Il y a en 3 types :
\begin{itemize}
 \item Plans \(\sigma_h\) : perpendiculaire à l'axe principal
 \item \(\sigma_v\) : contient l'axe principal
 \item \(\sigma_d\) : contient l'axe principal mais sont bissecteurs de 2 axes \(C_2'\)\marginCritical{Si il n'y a pas d'axes C2', pas de ce type de plan}
\end{itemize}
Les centres d'inversion sont le pendant des réflexion sauf que l'élément de symétrie est un point.
\subsection{Inversion}
L'élément est un centre d'inversion, un point.
\subsection{Rotation impropre}
\begin{definition}Elles se décomposent en 2 opérations :
\begin{enumerate}
 \item Une rotation \(C_n^m\) engendrée par un axe \(C_n\)
 \item Une réflexion par rapport à un plan \(\sigma\) orthogonal à l'axe \(C_n\)
\end{enumerate}
\(\sigma, C_n\) ne sont pas obligatoirement des éléments de symétrie de la molécule.
\end{definition}
%TODO Faire le schema pour le tétraèdre
\section{GPS}
L'ensemble des opérations de Symétrie des molécule constitue un groupe.

On détermine le GPS de chaque molécule avec l'organigramme de détermination des GPS.

Exemple : \begin{center}
\begin{tabular}{ll}
Molécule & GPS\\
H2O & C2V\\
NH3 & C3V\\
BH3 & D3H\\
Benzène & D6H\\
Diborane (B2H6) & D2H\\
CH4 & Td\\
PH5 & D3H\\
SH6 & Oh
          \end{tabular}
          \end{center}

\section{Application aux DOM}
\subsection{Représentation matricielle}
\begin{definition}[Représentation matricielle]
C'est l'ensemble des matrices qui représentent les opérations de symétrie du GPS. Il est noté \(\Gamma_n\) avec \(n\) la dimension du problème
\end{definition}
\(n\) correspond donc à la dimension des matrices carrées qui composent \(\Gamma\)
\subsubsection{3 dimensions}
Par exemple, la matrice de \(C_2\) selon l'axe z est \( \begin{pmatrix}
-1 & 0 & 0 \\
0 & -1 & 0 \\
0 & 0 & 1 \\
\end{pmatrix}\)
La présence de 1 signifie qu'il n'y a pas de transformions et le \(\-1\) montre que l'on a changé de signe pour la coordonnée sur l'axe. On peut réécrire sous la forme de matrice colonne.
\( \begin{pmatrix}
-1 \\
-1 \\
1 \\
\end{pmatrix}\).

la matrice de \(\sigma_v((xz))\) est \( \begin{pmatrix}
1 & 0 & 0 \\
0 & -1 & 0 \\
0 & 0 & 1 \\
\end{pmatrix}\)
\subsubsection{6 dimensions}
On utilise la base des OA de valence pur obtenir des matrices de dimension 6, dans la base \(\{2s,2p_x,2p_y,2p_z,1s_1,1s_2\}\). Dans cette base, \(C_2\) s'écrit  \( \begin{pmatrix}
1 & 0 & 0 &0&0&0 \\
0 & -1 & 0&0&0&0 \\
0 & 0 &- 1 &0&0&0\\
0&0&0&0&0&1\\
0&0&0&0&1&0
\end{pmatrix}\).
On regarde ce que fait la transformation pour chaque OA.

On remarque qu'il y a des éléments diagonaux nuls. On va faire une combinaison linéaire des OA qui s'échangent pour enlever les 0.

\subsection{Caractère de la représentation et RI}
\subsubsection{RI}
\begin{definition}[RI]
Chaque n-uplet est une représentation irréductible (RI) du groupe, irréductible car ne peut être
décomposé en matrices de dimensions plus petites. Ce sont les OS appliquées aux vecteur de la base.
\end{definition}
Par exemple, pour \(M(E) = \begin{pmatrix}
1\\1\\1
\end{pmatrix}, M(C_2) = \begin{pmatrix}
-1\\-1\\1
\end{pmatrix},  M(\sigma_v(xz)) = \begin{pmatrix}
1\\-1\\1
\end{pmatrix}, M(\sigma_v(yz)) = \begin{pmatrix}
-1\\1\\1
\end{pmatrix}\),

on peut écrire les quadruplets comme \(\begin{bmatrix}
[+1;-1;+1;-1;]\\
[+1;-1;-1;+1]\\
[+1;+1;+1;+1]
\end{bmatrix}\)

Chaque RI possède une dénomination propre et décrit comment un élément de la base se transforme selon les OS du GPS.

Chaque GPS possède un nombre fini de RI dont le nombre d’éléments dépend du nombre d’opérations de
symétrie du GPS.

\warningInfo{Représentation matricielle et RI}{La représentation matricielle est une représentation réductible (RR) car elle peut être décomposée
en une combinaison linéaire à coefficients entiers de RI.}

\subsubsection{Caractère}
C'est la trace de l'opération dans la base considérée. Ainsi, pour \(M(E), M(C_2), M(\sigma_v(xz)), M(\sigma_v(yz))\), le caractère \(\Gamma\) s'écrit \(\Gamma = \{3,-1,1,1\}\).
\subsection{Construction d'un DOM}
\subsubsection{Création du tableau}
À l'aide de la table de caractère du GPS de la molécule, on crée un tableau pour renseigner le résultat des opérations sur les OA mises en jeu dans la molécule.

Ainsi, pour la molécule d'H2O, on a
\begin{center}
\begin{tabular}{lllll}
 & E & C2 & \(\sigma_v(xz)\) & \(\sigma_v(yz)\)\\
2s & 1 & 1 & 1 & 1\\
2px & 1 & -1 & 1 & 1\\
2py & 1 & -1 & -1 & 1\\
2pz & 1 & 1 & 1 & 1\\
1s1 & 1 & 0 & 0 & 1\\
1s2 & 1 & 0 & 0 & 1
\end{tabular}
\end{center}
\warningInfo{Astuces}{l'opération E produit toujours des 1 et l'OA 2s centrale également car elle possède une symétrie sphérique.}
On calcule ensuite le caractère \(\Gamma\) pour chaque RI : Ici, on obtient \[\Gamma = \{6, 0, 2, 4\}\]
\begin{theorem}[Décomposition du Caractère]
\(n_i = \frac{1}{g}\sum_{k=1}^g \chi(R_K)\chi_i(R_k)\)
\end{theorem}
avec g le nombre total d'OS ici 4 car il y a (E, C2, \(\sigma_v(xz)\), \(\sigma_v(yz)\)) et  \(R_K\) une OS.

On utilise la table de caractère et le tableau construit.

Exemple : \(n_{A_1} = \frac{1}{4}(1\times1\times 6 + 1\times1\times 0 + 1\times1\times 2 + 1\times1\times 4) = 3\)

Pour \(n_{A_2} = \frac{1}{4}(1\times1\times 6 + 1\times1\times 0 + -1\times1\times 2 + -1\times1\times 4) = 0\)

\(n_{B1} = \frac{1}{4}(1\times1\times 6 + -1\times1\times 0 + 1\times1\times 2 + -1\times1\times 4) = 2\)

\(n_{B2} = \frac{1}{4}(1\times1\times 6 + -1\times1\times 0 - 1\times1\times 2 + 1\times1\times 4) = 1\)

Donc on peut écrire \(\Gamma_{RI} = 3A_1\oplus 0A_2\oplus 2B_1\oplus 1B_2\) La somme des résultat fait bien la dimension du problème.

On cherche maintenant l'emplacement de x,y, et z en avant dernière colonne du tableau pour déterminer. Ainsi, pour z dans la ligne \(A_1\), on aura \(2p_x = B1\), etc.

On modifie le tableau avec des changements de base pour retrouver le \(\Gamma_{RI}\). On fait des combinaisons linéaires pour trouver le nouveau tableau diagonalisé.
 \end{document}
